=============================================================================
% APPENDIX
% =============================================================================
\newpage
\appendix

\section{Training Data Generation Details}
\label{app:training}

Training data is generated by sampling 50,000 physically valid pairwise configurations
from \texttt{Rebound} \citep{Rein2012}. For each configuration, reference pairwise
forces are computed analytically using unsoftened Newtonian gravity. The configurations
are drawn from two distance regimes:

\begin{itemize}
  \item \textbf{Close-range ($60\%$)}: $r \in [0.5\epsilon,\, 50\epsilon]$
        where $\epsilon = 3\times10^{-4}$\,AU. In this regime the correction factor
        $c = (r_{\mathrm{soft}}/r)^3$ deviates meaningfully from unity, providing
        a non-trivial learning signal.
  \item \textbf{Far-field ($40\%$)}: $r \in [50\epsilon,\, 10]$\,AU. In this regime
        $c \approx 1$, and the NN is expected to learn near-unity outputs.
\end{itemize}

For each configuration, the target output is $\log c$ where $c$ is the ratio of the
true softened-to-unsoftened force magnitudes. The loss function is
$\mathcal{L}(\theta) = \mathrm{MSE}(\log \hat{c},\, \log c_{\mathrm{true}})$.

The training convergence curve is shown in Figure~\ref{fig:training_log}.

\begin{figure}[H]
    \centering
    \includegraphics[width=0.75\linewidth]{training_log.png}
    \caption{Training and validation loss curves over 5,000 epochs for seed~42 (the
    model used throughout this paper). Both curves converge smoothly to
    $\approx 10^{-5}$ without instability or overfitting. Best validation MSE:
    $1.11\times10^{-5}$.}
    \label{fig:training_log}
\end{figure}

\section{Supplementary Figures}
\label{app:figures}

Figure~\ref{fig:ic4_timeseries} shows the energy drift time series for IC4 across
all five tested timesteps, demonstrating the $dt^2$ scaling relationship and the
integration breakdown at $dt = 0.08$\,yr.

\begin{figure}[H]
    \centering
    \includegraphics[width=0.95\linewidth]{02_drift_timeseries_IC4_hierarchical.png}
    \caption{IC4 (hierarchical) energy drift $|\Delta E / E_0|$ over 100 years at
    all five tested timesteps. $dt = 0.01$\,yr (green) and $dt = 0.04$\,yr (orange,
    operational) show smooth oscillating leapfrog drift. $dt = 0.08$\,yr and
    $dt = 0.10$\,yr (red, purple dashed) show a permanent $\sim 50\%$ energy offset
    from $t \approx 10$\,yr, indicating genuine integration breakdown at those
    timesteps. This confirms that the chosen operational timestep $dt = 0.04$\,yr remains viable for this configuration, while the larger tested timesteps ($dt = 0.08$ and $0.10$\,yr) break down.}
    \label{fig:ic4_timeseries}
\end{figure}

second try
% ============================================================
% APPENDIX A: Initial Condition Details
% ============================================================
\appendix
\section{Initial Condition Details}
\label{app:ics}

Table~\ref{tab:ic_details} provides the full specification of all seven initial conditions
evaluated in Section~\ref{sec:generalisation}. All coordinates are after centre-of-mass centering
and are given in astronomical units (AU) and AU\,yr$^{-1}$.

\begin{table}[H]
\centering
\caption{Full initial condition specifications. All masses in $\mathrm{M}_\odot$; coordinates in
AU and AU/yr after CoM centering. $E_0$: total initial energy in natural units ($G=1$).}
\label{tab:ic_details}
\scriptsize
\setlength{\tabcolsep}{4pt}
\renewcommand{\arraystretch}{1.12}

\begin{tabularx}{\textwidth}{@{}l l X r l@{}}
\toprule
\textbf{IC} & \textbf{Masses} & \textbf{Configuration} & \textbf{$E_0$} & \textbf{Outcome} \\
\midrule
IC1 & [1.0, 0.01, 0.005] &
$x_0$=[[0,0,0],[1,0,0],[0,1.2,0]], $v_0$=[[0,0,0],[0,1,0],[$-$0.9,0,0]] &
$-$0.00723 & Bounded \\

IC2 & [1.0, 0.5, 0.25] &
$x_0$=[[0,0,0],[1,0,0],[$-$0.5,0.8,0]], $v_0$=[[0,0,0],[0,0.6,0],[$-$0.4,$-$0.3,0]] &
$-$0.73460 & Unstable \\

IC3 & [1.0, 0.01, 0.005] &
$x_0$=[[0,0,0],[0.5,0,0],[0,2.5,0]], $v_0$=[[0,0,0],[0,1.3,0],[$-$0.4,0,0]] &
$-$0.01325 & Bounded \\

IC4 & [1.0, 0.01, 0.005] &
$x_0$=[[0,0,0],[1,0,0],[0,5,0]], $v_0$=[[0,0,0],[0,1,0],[$-$0.12,0,0]] &
$-$0.00602 & Bounded \\

IC5 & [1.0, 0.01, 0.005] &
$x_0$=[[0,0,0],[0.3,0,0],[0,2,0]], $v_0$=[[0,0,0],[0,2.2,0],[$-$0.3,0,0]] &
$-$0.01167 & Unstable \\

IC6 & [1.0, 0.01, 0.005] &
$x_0$=[[0,0,0],[1.5,0,0],[$-$3,0,0]], $v_0$=[[0,0,0],[0,0.816,0],[0,$-$0.577,0]] &
$-$0.00420 & Bounded \\

IC7 & [1.0, 0.01, 0.005] &
$x_0$=[[0,0,0],[0.6,0,0],[0,4,0]], $v_0$=[[0,0,0],[0,1.6,0],[$-$0.15,0,0]] &
$-$0.00520 & Unstable \\
\bottomrule
\end{tabularx}
\end{table}

\section{Supplementary Figures}
\label{app:figs}

\begin{figure}[H]
    \centering
    \includegraphics[width=0.9\linewidth]{traj_overlay_timeseries_all_bodies_T100_after_adaptive.png}
    \caption{Coordinate time series overlays $x(t)$, $y(t)$ for all three bodies over
    100 years: \texttt{ias15} (solid) vs SIMON (dashed). Early-time agreement is strong;
    a phase offset develops after $\sim$25 years, consistent with Lyapunov-time chaotic
    divergence. All coordinates remain oscillatory and bounded throughout.}
    \label{fig:timeseries_app}
\end{figure}

\begin{figure}[H]
    \centering
    \includegraphics[width=0.9\linewidth]{timeseries_no_nn.png}
    \caption{Coordinate time series for SIMON (green dashed) vs No-NN baseline
    (red dotted) vs \texttt{ias15} (blue solid), all three bodies, $T=100$\,yr.
    SIMON and No-NN follow the same qualitative trajectory; the 8.5\% RMS difference
    manifests as a slightly larger phase offset in the No-NN case at later times.}
    \label{fig:timeseries_no_nn_app}
\end{figure}

third try
% ============================================================
% APPENDIX
% ============================================================
\appendix

\section{Initial Condition Coordinates}
\label{app:ic-details}

Table~\ref{tab:ic_coords} provides the full initial conditions for all seven configurations used in Section~\ref{sec:generalisation}.

\begin{table}[H]
\centering
\caption{Initial conditions for all seven configurations. Masses in $\mathrm{M}_\odot$; positions in AU; velocities in AU/yr. $E_{\rm total}$ in simulation units ($G=1$).}
\label{tab:ic_coords}
\small
\setlength{\tabcolsep}{4pt}
\renewcommand{\arraystretch}{1.12}

\begin{tabularx}{\textwidth}{@{}l l X r X@{}}
\toprule
\textbf{IC} & \textbf{Masses ($\mathrm{M}_\odot$)} & \textbf{Configuration} & \textbf{$E_{\rm total}$} & \textbf{Regime} \\
\midrule
IC1 (Default) &
$[1.0, 0.01, 0.005]$ &
Bodies at $[0,0]$, $[1,0]$, $[-1,0]$ &
$-0.00723$ &
Moderate scattering \\

IC2 (Near-equal) &
$[1.0, 0.5, 0.25]$ &
Bodies at $[0,0]$, $[1,0]$, $[-1,0]$ &
$-0.73460$ &
Strongly interacting \\

IC3 (Tight pair) &
$[1.0, 0.01, 0.005]$ &
Bodies at $[0,0]$, $[0.5,0]$, $[-1,0]$ &
$-0.01325$ &
Tight inner binary \\

IC4 (Hierarchical) &
$[1.0, 0.01, 0.005]$ &
Bodies at $[0,0]$, $[1,0]$, $[5,0]$ &
$-0.00602$ &
Outer body at 5\,AU \\

IC5 (High ecc) &
$[1.0, 0.01, 0.005]$ &
Body 1 at $[0.3,0]$, $v=2.2$ (85\% $v_{\rm esc}$) &
$-0.01167$ &
High eccentricity \\

IC6 (Circular) &
$[1.0, 0.01, 0.005]$ &
Keplerian circular orbits &
$-0.00420$ &
Near-circular \\

IC7 (High ecc v2) &
$[1.0, 0.01, 0.005]$ &
Body 1 at $[0.6,0]$, $v=1.6$ ($1.24\times v_{\rm circ}$) &
$-0.00520$ &
High ecc redesign \\
\bottomrule
\end{tabularx}
\end{table}

\section{Supplementary Figures}
\label{app:figures1}

\begin{figure}[H]
    \centering
    \includegraphics[width=0.9\linewidth]{multi_ic_trajectories_v3.png}
    \caption{Trajectory overlays (\textsc{simon} dashed, \texttt{ias15} solid)
    for all seven initial conditions. Bounded configurations (IC1, IC3, IC4, IC6)
    show qualitatively consistent bounded orbits. Unbounded configurations
    (IC2, IC5, IC7) show clear ejection trajectories.}
    \label{fig:app_traj}
\end{figure}

\begin{figure}[H]
    \centering
    \includegraphics[width=0.82\linewidth]{02_drift_timeseries_IC4_hierarchical.png}
    \caption{IC4 energy drift timeseries at five tested timestep values. At $dt=0.04$\,yr (operational, orange), drift reaches $13.77\%$ with oscillatory character (shadow Hamiltonian conserved). At $dt=0.01$\,yr (green), drift reduces to $1.05\%$. At $dt=0.08$\,yr and $dt=0.10$\,yr, integration breakdown occurs (permanent offset from $t\approx10$\,yr), showing that the chosen operational timestep $dt=0.04$\,yr remains viable for IC4, while larger tested timesteps break down.}
    \label{fig:app_ic4}
\end{figure}
